% !PRINTED 600
\weberpage{600}
Wenn nun

\[
\mathfrak{p}_{r,s} = \mathfrak{p}_r\,|\,\varphi_{r,s}
\]

gesetzt wird, so ist nach \S.\,160, (6)

\weq{6}{\mathfrak{p}_r' = (\mathfrak{p}_{r,1}\,\mathfrak{p}_{r,2}\cdots\mathfrak{p}_{r,e_r})^{g_r},}

und die Substitution von (6) in (3) ergiebt durch Vergleichung
mit \S.\,160, (6)

\weq{7}{g = a_r\,g_r, \quad e_1 + e_2 + \cdots + e_{e'} = e,}

wodurch die Exponenten $a_r$ definirt sind.

Die in Bezug auf den Körper $\Omega'$ genommene Partialnorm
von $\mathfrak{p}_{r,s}$ ist nach \S.\,160, (3), (7):

\weq{8}{\mathfrak{N}_{\Omega'}(\mathfrak{p}_{r,s}) = \mathfrak{p}_r'^{\,f_r},}

wenn $f_r$ durch die Gleichung

\weq{9}{n' = e_r\,f_r\,g_r, \quad (h_r = f_r\,g_r)}

definirt wird.

Wir wollen nun die Gruppe $\Phi$ nach $\Phi'$ in Nebengruppen zerlegen,
und setzen, wenn $\vartheta_1, \vartheta_2, \ldots, \vartheta_{m'}$ passend ausgewählte
Substitutionen aus $\Phi$ sind:

\weq{10}{\Phi = \Phi'\vartheta_1 + \Phi'\vartheta_2 + \cdots + \Phi'\vartheta_{m'},}

so dass durch die Substitutionen $\vartheta_1, \vartheta_2, \ldots, \vartheta_{m'}$ der Körper
$\Omega'$ in $m'$ conjugirte (gleiche oder verschiedene) Körper

\[
\Omega_1', \ \Omega_2', \ \ldots, \ \Omega_{m'}'
\]

übergeht. Der Körper $\Omega_t'$ gehört dann zu der Gruppe $\vartheta_t^{-1}\Phi'\vartheta_t$.

Wenn nun durch $\vartheta_t$ die Ideale $\mathfrak{p}_r'$ und $\mathfrak{p}_{r,s}$ in $\mathfrak{p}_{r,t}'$ und
$\mathfrak{p}_{r,s,t}$ übergehen, so ist nach \S.\,160, (6)

\weq{11}{\mathfrak{p}_{r,t}' = (\mathfrak{p}_{r,1,t}\,\mathfrak{p}_{r,2,t}\cdots\mathfrak{p}_{r,e_r,t})^{g_r},}

und das Product aller dieser Ideale für $t = 1, 2, \ldots, m'$ ist die
im Körper $\Omega'$ genommene Partialnorm von $\mathfrak{p}_r'$ in Bezug auf den
Körper $R$, die wir mit $\mathfrak{N}_R'(\mathfrak{p}_r')$ bezeichnen. Sie muss eine Potenz
von $\mathfrak{P}$ sein, und wir setzen

\weq{12}{\mathfrak{N}_R'(\mathfrak{p}_r') = \mathfrak{P}^{\,f_r'}.}

Um $f_r'$ zu finden, brauchen wir nur $\mathfrak{P}$ in (12) nach \S.\,160, (6)
in seine Primfactoren $\mathfrak{p}$ zu zerlegen, wodurch sich $e\,g\,f_r'$ Primfactoren
 $\mathfrak{p}$ in $\mathfrak{P}^{\,f_r'}$ ergeben. Bilden wir andererseits das Product
der $m'$ Factoren (11), so erhalten wir $g_r\,e_r\,m'$ solcher Primfactoren.
Es ist daher

\[
e\,g\,f_r' = m'\,g_r\,e_r,
\]

also nach (1), (9) und \S.\,160, (7):

\weq{13}{f = f_r\,f_r'.}

\begin{center}
\rule{0.3\linewidth}{0.4pt}
\end{center}
